% Source: https://artofproblemsolving.com/community/c3h2251656_a_problem_i_made_up
% Stable identity: aops-topic-2251656-q01
Do (b) first, with side $2x$. Put the square with corners $(0,0)$ and $(2x,2x)$; the big circle has centre $O=(x,x)$ and radius $x$. The small circle in the corner at the origin touches both axes, so its centre is $(r,r)$, where $r$ is its radius; it lies on the diagonal from the corner to $O$.

\begin{asy}
usepackage("amsmath");
size(5cm);
defaultpen(fontsize(10pt));
real s=10, R=5, r=5*(3-2*sqrt(2));
pair O=(R,R);
draw((0,0)--(s,0)--(s,s)--(0,s)--cycle,linewidth(0.9));
draw(circle(O,R),linewidth(0.9));
pair[] cs={(r,r),(s-r,r),(s-r,s-r),(r,s-r)};
for(pair c:cs) filldraw(circle(c,r),palecyan,blue+linewidth(0.8));
draw((0,0)--O,red+dashed);
dot(O); dot((r,r),blue); label("$O$",O,NE);
label("$10$",(s/2,s),N);
draw((r,r)--(r,0),blue); label("$r$",(r,r/2),E,blue+fontsize(8pt));
\end{asy}

The two circles are tangent externally, so the distance between their centres is $x+r$:
\[\sqrt2\,(x-r)=x+r\quad\Longrightarrow\quad r=\dfrac{\sqrt2-1}{\sqrt2+1}\,x=(\sqrt2-1)^2x=(3-2\sqrt2)\,x.\]
Then $r^2=(17-12\sqrt2)\,x^2$, and the area inside the square but outside the four small circles is
\[(2x)^2-4\pi r^2=\boxed{4x^2-4\pi(17-12\sqrt2)\,x^2}=4x^2\bigl(1-(17-12\sqrt2)\pi\bigr).\]
This answers \textbf{(b)}. For \textbf{(a)}, $x=5$:
\[100-100\pi(17-12\sqrt2)=\boxed{100-(1700-1200\sqrt2)\pi}\approx90.75.\]
(Each small circle has radius $5(3-2\sqrt2)\approx0.858$.)
